\documentclass[../../数学分析培优讲义.tex]{subfiles}

\begin{document}

\setcounter{chapter}{5}
\pagestyle{fancy}

\chapter{定积分}


\section{积分计算}


\subsection{对称性}

\begin{proposition}[\marktwostar]
记
\[
    I(m,n)=\int_0^{\frac{\pi}{2}}\cos^m x\sin^n x\,\mathrm{d}x,
\]
则 $I(m,n)=I(n,m)$,且
\[
    I(m,n)=\dfrac{m-1}{m+n}I(m-2,n)
    =\dfrac{n-1}{m+n}I(m,n-2),
\]
并有
\[
I(m,n)=
\begin{cases}
\dfrac{(m-1)!!(n-1)!!}{(m+n)!!},&m,n\text{ 不全为偶数},\\[6pt]
\dfrac{(m-1)!!(n-1)!!}{(m+n)!!}\dfrac{\pi}{2},&m,n\text{ 都为偶数}.
\end{cases}
\]
\end{proposition}

\begin{theorem}[\markstar]
设 $f(x),g(x)$ 为两个实函数,那么:
\begin{enumerate}[label=(\arabic*)]
    \item $f(x)$ 与 $g(x)$ 都关于 $x=x_0$ 轴对称,则 $f(x)g(x)$ 也关于 $x=x_0$ 轴对称;
    \item $f(x)$ 关于 $x=x_0$ 轴对称,$g(x)$ 关于 $(x_0,0)$ 中心对称,则 $f(x)g(x)$ 关于 $(x_0,0)$ 中心对称;
    \item $f(x)$ 与 $g(x)$ 都关于 $(x_0,0)$ 中心对称,则 $f(x)g(x)$ 关于 $x=x_0$ 轴对称.
\end{enumerate}
\end{theorem}

\begin{corollary}[\markstar]
对任意非负整数 $m,n$,有:
\begin{enumerate}[label=(\arabic*)]
    \item $\sin^n x$ 关于 $x=\dfrac{\pi}{2}$ 轴对称;
    \item $m$ 为偶数时,$\cos^m x$ 关于 $x=\dfrac{\pi}{2}$ 轴对称;$m$ 为奇数时,$\cos^m x$ 关于 $\left(\dfrac{\pi}{2},0\right)$ 中心对称;
    \item $m$ 为偶数时,$\cos^m x\sin^n x$ 关于 $x=\dfrac{\pi}{2}$ 轴对称;当 $m$ 为奇数时,$\cos^m x\sin^n x$ 关于 $\left(\dfrac{\pi}{2},0\right)$ 中心对称.
\end{enumerate}
\end{corollary}

\begin{proposition}[\markstar]
对于任意非负整数 $m,n$,有
\[
    J(m,n)\overset{\triangle}{=}
    \int_0^\pi\cos^m x\sin^n x\,\mathrm{d}x
    =\begin{cases}
        0,&m\text{ 为奇数},\\
        2I(m,n),&m\text{ 为偶数}.
    \end{cases}
\]
\end{proposition}

\begin{remark}
此时关于 $m,n$ 并非对称,即一般情况下 $J(m,n)\neq J(n,m)$.
\end{remark}

\begin{proposition}[\markstar]
对于任意非负整数 $m,n$,有
\[
    K(m,n)\overset{\triangle}{=}
    \int_0^{2\pi}\cos^m x\sin^n x\,\mathrm{d}x
    =\begin{cases}
        0,&n\text{ 为奇数},\\
        2J(m,n),&n\text{ 为偶数},
    \end{cases}
\]
从而
\[
    K(m,n)=\begin{cases}
        0,&m,n\text{ 不全为偶数},\\
        4I(m,n),&m,n\text{ 都为偶数}.
    \end{cases}
\]
\end{proposition}

\begin{example}
计算积分:
\begin{enumerate}[label=(\arabic*)]
    \item $\displaystyle\int_0^{\frac{\pi}{2}}\cos^3x\sin^2x\,\mathrm{d}x$;
    \item $\displaystyle\int_0^{\frac{\pi}{2}}\cos^4x\sin^3x\,\mathrm{d}x$;
    \item $\displaystyle\int_0^{\frac{\pi}{2}}\sin^3x\cos^5x\,\mathrm{d}x$;
    \item $\displaystyle\int_0^{\frac{\pi}{2}}\sin^2x\cos^4x\,\mathrm{d}x$.
\end{enumerate}
\end{example}
\exampleblank

\begin{example}
计算积分:
\begin{enumerate}[label=(\arabic*)]
    \item $\displaystyle\int_0^\pi\cos^3x\sin^2x\,\mathrm{d}x$;
    \item $\displaystyle\int_0^\pi\cos^4x\sin^3x\,\mathrm{d}x$;
    \item $\displaystyle\int_0^\pi\cos^n x\,\mathrm{d}x$;
    \item $\displaystyle\int_0^\pi\sin^n x\,\mathrm{d}x$.
\end{enumerate}
\end{example}
\exampleblank

\begin{example}
计算积分:
\begin{enumerate}[label=(\arabic*)]
    \item $\displaystyle\int_0^{2\pi}\cos^4x\sin^3x\,\mathrm{d}x$;
    \item $\displaystyle\int_0^{2\pi}\sin^3x\cos^5x\,\mathrm{d}x$;
    \item $\displaystyle\int_0^{2\pi}\sin^2x\cos^4x\,\mathrm{d}x$;
    \item $\displaystyle\int_0^{2\pi}\sin^n x\,\mathrm{d}x
    =\int_0^{2\pi}\cos^n x\,\mathrm{d}x$.
\end{enumerate}
\end{example}
\exampleblank

\begin{example}
计算积分:
\begin{enumerate}[label=(\arabic*)]
    \item $\displaystyle\int_{-\pi}^{\pi}\cos^4x\sin^4x\,\mathrm{d}x$;
    \item $\displaystyle\int_{-\pi}^{3\pi}\sin^5x\cos^6x\,\mathrm{d}x$;
    \item $\displaystyle\int_{-\pi}^{0}\cos^3x\sin^2x\,\mathrm{d}x$;
    \item $\displaystyle\int_0^\pi|\cos^5x|\,\mathrm{d}x$.
\end{enumerate}
\end{example}
\exampleblank

\begin{example}
计算积分:
\begin{enumerate}[label=(\arabic*)]
    \item $\displaystyle\int_0^\pi\dfrac{\mathrm{d}x}{1+\sin^2x}$;
    \item $\displaystyle\int_0^{2\pi}\dfrac{\mathrm{d}x}{1+a\cos x}\quad(-1<a<1)$.
\end{enumerate}
\end{example}
\exampleblank

\begin{example}
利用对称性,计算下列积分:
\begin{enumerate}[label=(\arabic*)]
    \item $\displaystyle\int_0^\pi\dfrac{x}{1+\cos^2x}\,\mathrm{d}x$;
    \item $\displaystyle\int_0^1\dfrac{x}{e^x+e^{1-x}}\,\mathrm{d}x$;
    \item $\displaystyle\int_{-2}^{2}x\ln(1+e^x)\,\mathrm{d}x$;
    \item $\displaystyle\int_0^{\frac{\pi}{4}}\ln(1+\tan x)\,\mathrm{d}x$;
    \item $\displaystyle\int_0^{\frac{\pi}{2}}\dfrac{\sin^n x}{\sin^n x+\cos^n x}\,\mathrm{d}x$;
    \item $\displaystyle\int_0^\pi
    \dfrac{a^n\sin^2x+b^n\cos^2x}{a^{2n}\sin^2x+b^{2n}\cos^2x}\,\mathrm{d}x$.
\end{enumerate}
\end{example}
\exampleblank

\begin{example}
设 $f\in C(\lbrack 0,a\rbrack)$,$a>0$.
\begin{enumerate}[label=(\arabic*)]
    \item 在 $\lbrack 0,a\rbrack$ 上 $f(x)+f(a-x)\neq0$,计算
    \[
        I=\int_0^a\dfrac{f(x)}{f(x)+f(a-x)}\,\mathrm{d}x;
    \]
    \item 在 $\lbrack 0,a\rbrack$ 上 $f(x)f(a-x)\equiv1$,计算
    \[
        I=\int_0^a\dfrac{\mathrm{d}x}{1+f(x)}.
    \]
\end{enumerate}
\end{example}
\exampleblank

\begin{example}
设 $f$ 为连续函数,证明下列等式:
\begin{enumerate}[label=(\arabic*)]
    \item
    \[
        \int_1^a f\left(x^2+\dfrac{a^2}{x^2}\right)\dfrac{\mathrm{d}x}{x}
        =\int_1^a f\left(x+\dfrac{a^2}{x}\right)\dfrac{\mathrm{d}x}{x};
    \]
    \item
    \[
        \int_0^{2\pi}f(a\cos x+b\sin x)\,\mathrm{d}x
        =2\int_0^\pi f\left(\sqrt{a^2+b^2}\cos x\right)\,\mathrm{d}x.
    \]
\end{enumerate}
\end{example}
\exampleblank



\subsection{积分与极限}

\begin{example}
设 $f(x)$ 在 $\lbrack 0,1\rbrack$ 上连续,求证:
\[
    \lim\limits_{n\to\infty}\dfrac{1}{n}
    \sum\limits_{k=1}^{n-1}(-1)^{k+1}f\left(\dfrac{k}{n}\right)=0.
\]
\end{example}
\exampleblank

\begin{example}
求极限
\[
    \lim\limits_{n\to\infty}
    \left(b^{\frac{1}{n}}-1\right)
    \sum\limits_{i=0}^{n-1}
    b^{\frac{1}{n}}\sin b^{\frac{1}{n}}
    \sin b^{\frac{n+1}{2n}}.
\]
\end{example}
\exampleblank

\begin{example}
求极限
\[
    \lim\limits_{x\to+\infty}\dfrac{1}{x}\int_0^x|\sin t|\,\mathrm{d}t.
\]
\end{example}
\exampleblank

\begin{example}
求证:
\begin{enumerate}[label=(\arabic*)]
    \item $\displaystyle\lim\limits_{n\to\infty}\int_0^{\frac{\pi}{2}}\sin^n x\,\mathrm{d}x=0$;
    \item $\displaystyle\lim\limits_{n\to\infty}\int_0^{\frac{\pi}{2}}(1-\sin x)^n\,\mathrm{d}x=0$,其中 $a\in(0,1)$;
    \item $\displaystyle\lim\limits_{n\to\infty}\int_0^1e^{x^n}\,\mathrm{d}x=1$.
\end{enumerate}
\end{example}
\exampleblank

\begin{example}
设 $f(x)$ 严格单调递减,在 $\lbrack 0,1\rbrack$ 上连续,$f(0)=1,f(1)=0$.试证明:对任意 $\delta\in(0,1)$ 有
\begin{enumerate}[label=(\arabic*)]
    \item
    \[
        \lim\limits_{n\to\infty}
        \dfrac{\displaystyle\int_\delta^1[f(x)]^n\,\mathrm{d}x}
        {\displaystyle\int_0^\delta[f(x)]^n\,\mathrm{d}x}=0;
    \]
    \item
    \[
        \lim\limits_{n\to\infty}
        \dfrac{\displaystyle\int_0^\delta[f(x)]^{n+1}\,\mathrm{d}x}
        {\displaystyle\int_0^1[f(x)]^n\,\mathrm{d}x}=1.
    \]
\end{enumerate}
\end{example}
\exampleblank

\begin{example}[\markstar]
设 $f(x)\geq0,g(x)>0$,两函数在 $\lbrack a,b\rbrack$ 上连续,求证:
\[
    \lim\limits_{n\to\infty}
    \left[\int_a^b[f(x)]^n g(x)\,\mathrm{d}x\right]^{\frac{1}{n}}
    =\max\limits_{a\leq x\leq b}f(x).
\]
\end{example}
\exampleblank

拟合法:当证明极限 $\displaystyle\lim\limits_{n\to\infty}a_n=a$ 时,若将 $a$ 写成类似 $a_n$ 的形式,则通过证明 $\displaystyle\lim\limits_{n\to\infty}(a_n-a)=0$,将会大大简化问题.

\begin{example}
设 $f(x)$ 在 $\lbrack 0,1\rbrack$ 上连续,试证:
\[
    \lim\limits_{h\to0^+}\int_0^1\dfrac{h}{h^2+x^2}f(x)\,\mathrm{d}x
    =\dfrac{\pi}{2}f(0).
\]
\end{example}
\exampleblank

\begin{example}
设 $f$ 为连续函数,证明:
\[
    \lim\limits_{n\to\infty}\dfrac{2}{\pi}
    \int_0^1\dfrac{n}{n^2x^2+1}f(x)\,\mathrm{d}x=f(0).
\]
\end{example}
\exampleblank

\begin{example}
设 $f(x)$ 在 $\lbrack a,b\rbrack$ 上可积,在 $x=b$ 处左连续,求证:
\[
    \lim\limits_{n\to\infty}
    \dfrac{n+1}{(b-a)^{n+1}}
    \int_a^b(x-a)^n f(x)\,\mathrm{d}x=f(b).
\]
\end{example}
\exampleblank

\begin{example}
设 $f\in C\lbrack A,B\rbrack$,$A<a<b<B$,试证:
\[
    \lim\limits_{h\to0}\int_a^b\dfrac{f(x+h)-f(x)}{h}\,\mathrm{d}x
    =f(b)-f(a).
\]
\end{example}
\exampleblank

\begin{example}
设 $f(x)$ 在 $\lbrack a,b\rbrack$ 上可导,$f''(x)$ 在 $\lbrack a,b\rbrack$ 上可积.对任意 $n\in\mathbb{N}$,记
\[
    A_n=\sum\limits_{i=1}^n f\left(a+i\dfrac{b-a}{n}\right)\dfrac{b-a}{n}
    -\int_a^b f(x)\,\mathrm{d}x.
\]
试证:
\[
    \lim\limits_{n\to\infty}nA_n
    =\dfrac{b-a}{2}[f(b)-f(a)].
\]
\end{example}
\exampleblank

\begin{example}[\markstar]
设 $f(x)$ 在 $\lbrack a,b\rbrack$ 上二次可微,且 $f''(x)$ 在 $\lbrack a,b\rbrack$ 上可积,记
\[
    B_n=\int_a^b f(x)\,\mathrm{d}x
    -\dfrac{b-a}{n}\sum\limits_{i=1}^n
    f\left[a+(2i-1)\dfrac{b-a}{2n}\right].
\]
试证:
\[
    \lim\limits_{n\to\infty}n^2B_n
    =\dfrac{(b-a)^2}{24}[f'(b)-f'(a)].
\]
\end{example}
\exampleblank

\begin{remark}
因此,若设
\[
    A_n=\sum\limits_{k=1}^n\dfrac{1}{n+k},
    \qquad
    B_n=\sum\limits_{k=1}^n\dfrac{2}{2n+2k-1},
\]
则
\[
    \lim\limits_{n\to\infty}n(\ln2-A_n)=\dfrac{1}{4},
    \qquad
    \lim\limits_{n\to\infty}n^2(\ln2-B_n)=\dfrac{1}{32}.
\]
\end{remark}

\begin{example}[\markstar]
设 $f(x)$ 在 $\lbrack a,b\rbrack$ 上可积,$g(x)\geq0$,且 $g$ 是以 $T>0$ 为周期、在 $\lbrack 0,T\rbrack$ 上可积的函数.试证:
\[
    \lim\limits_{n\to\infty}\int_a^b f(x)g(nx)\,\mathrm{d}x
    =\dfrac{1}{T}\int_0^Tg(x)\,\mathrm{d}x\int_a^b f(x)\,\mathrm{d}x.
\]
\end{example}
\exampleblank

\begin{example}[\markstar\ Riemann 引理]
若 $f(x)$ 在 $\lbrack a,b\rbrack$ 上可积,$g(x)$ 以 $T$ 为周期,在 $\lbrack 0,T\rbrack$ 上可积,则
\[
    \lim\limits_{n\to\infty}\int_a^b f(x)g(nx)\,\mathrm{d}x
    =\dfrac{1}{T}\int_0^Tg(x)\,\mathrm{d}x\int_a^b f(x)\,\mathrm{d}x.
\]
\end{example}
\exampleblank

\begin{example}
设函数 $f(x)$ 在区间 $\lbrack 0,1\rbrack$ 上有一阶连续导数,且 $f(0)=f(1)$,$g(x)$ 是周期为 $1$ 的连续函数,记
\[
    a_n=\int_0^1 f(x)g(x)\,\mathrm{d}x,
\]
求证:
\[
    \lim\limits_{n\to\infty}na_n=0.
\]
\end{example}
\exampleblank

\begin{example}
设 $s(x)=4[x]-2[2x]+1$,$f$ 在 $\lbrack 0,1\rbrack$ 上可积,证明:
\[
    \lim\limits_{n\to\infty}\int_0^1 f(x)s(nx)\,\mathrm{d}x=0.
\]
\end{example}
\exampleblank


\newpage

\section{可积性}

\begin{theorem}[\markstar]
$\lbrack a,b\rbrack$ 上的有界函数可积的充分必要条件是:对任意 $\varepsilon>0,\sigma>0$,存在划分
\[
    \Delta:a=x_0<x_1<\cdots<x_n=b,
\]
其相应于 $\omega\geq\varepsilon$ 的子区间 $\Delta x_i$ 的长度总和小于 $\sigma$.
\end{theorem}

\begin{corollary}[\markstar]
若定义在 $\lbrack a,b\rbrack$ 上的有界函数只有有限个不连续点,则 $f(x)$ 在 $\lbrack a,b\rbrack$ 上可积.
\end{corollary}

\begin{example}
试证明下列函数 $f(x)$ 在区间 $\lbrack 0,1\rbrack$ 上可积:
\begin{enumerate}[label=(\arabic*)]
    \item
    \[
        f(x)=\begin{cases}
            \sin\dfrac{1}{x},&x\neq0,\\
            0,&x=0;
        \end{cases}
    \]
    \item
    \[
        f(x)=\begin{cases}
            \dfrac{1}{x}-\dfrac{1}{\sin x},&x\neq0,\\
            0,&x=0;
        \end{cases}
    \]
    \item
    \[
        f(x)=\begin{cases}
            \ln x\cdot\ln(1+x),&x\neq0,\\
            0,&x=0.
        \end{cases}
    \]
\end{enumerate}
\end{example}
\exampleblank

\begin{proposition}[\markstar]
设 $f(x)$ 为 $\lbrack a,b\rbrack$ 上的有界函数,$\{a_n\}\subseteq\lbrack a,b\rbrack$,且 $\displaystyle\lim\limits_{n\to\infty}a_n=c$.证明:若 $f(x)$ 在 $\lbrack a,b\rbrack$ 上只有 $a_n$($n=1,2,3,\ldots$)为其间断点,则 $f(x)$ 在 $\lbrack a,b\rbrack$ 上可积.
\end{proposition}

\begin{example}
试证明下列函数在 $\lbrack 0,1\rbrack$ 上是可积的:
\begin{enumerate}[label=(\arabic*)]
    \item
    \[
        f(x)=\begin{cases}
            0,&x=0,\\
            \dfrac{1}{x}-\left[\dfrac{1}{x}\right],&x\neq0;
        \end{cases}
    \]
    \item $f(x)=\operatorname{sgn}\left(\sin\dfrac{\pi}{x}\right)$.
\end{enumerate}
\end{example}
\exampleblank

\begin{example}[\markstar]
证明:Riemann 函数在 $\lbrack 0,1\rbrack$ 上可积.
\end{example}
\exampleblank

\begin{theorem}[\markstar\ 四则运算]
若 $f(x),g(x)$ 在 $\lbrack a,b\rbrack$ 上可积,则:
\begin{enumerate}[label=(\arabic*)]
    \item $kf(x),f(x)\pm g(x),f(x)g(x)$ 在 $\lbrack a,b\rbrack$ 上也可积,但 $\dfrac{f(x)}{g(x)}$($g(x)\neq0$)也不一定可积;
    \item 若 $|f(x)|\geq m>0$,则 $\dfrac{1}{f(x)}$ 在 $\lbrack a,b\rbrack$ 上也可积;
    \item 若 $f(x)$ 在 $\lbrack a,b\rbrack$ 上连续且恒不为零,则 $\dfrac{1}{f(x)}$ 在 $\lbrack a,b\rbrack$ 上也可积.
\end{enumerate}
\end{theorem}

\begin{example}[\markstar\ 复合]
设 $f(x)$ 在 $\lbrack a,b\rbrack$ 上连续,$\varphi(t)$ 在 $\lbrack \alpha,\beta\rbrack$ 上可积,且
\[
    a\leq\varphi(t)\leq b,
    \qquad t\in\lbrack \alpha,\beta\rbrack,
\]
则 $f(\varphi(t))$ 在 $\lbrack \alpha,\beta\rbrack$ 上仍然可积.
\end{example}
\exampleblank

\begin{remark}
可积函数的复合不一定是可积函数,例如 $f(x)=\operatorname{sgn}x$ 与 $\varphi(t)=R(t)$ 为 Riemann 函数.
\end{remark}

\begin{example}[\markstar]
讨论 $f(x),|f(x)|,f^2(x)$ 在 $\lbrack a,b\rbrack$ 上可积性的关系.
\end{example}
\exampleblank

\begin{proposition}[\markstar]
设 $f\in R\lbrack a,b\rbrack$,证明:$f$ 的连续点在 $\lbrack a,b\rbrack$ 中稠密.
\end{proposition}

\begin{corollary}[\markstar]
设非负函数 $f\in R\lbrack a,b\rbrack$,则
\[
    \int_a^b f(x)\,\mathrm{d}x=0
    \quad\Longleftrightarrow\quad
    f\text{ 在连续点上恒为零}.
\]
\end{corollary}

\begin{example}
设 $f(x)$ 在 $\lbrack a,b\rbrack$ 上可积,试证:对于 $\lbrack a,b\rbrack$ 上任意可积函数 $g(x)$,恒有
\[
    \int_a^b g(x)f(x)\,\mathrm{d}x=0,
\]
则函数 $f(x)$ 在连续点上恒为零.
\end{example}
\exampleblank

\begin{example}
设在 $\lbrack -1,1\rbrack$ 上的连续函数 $f(x)$ 满足如下条件:对 $\lbrack -1,1\rbrack$ 上的任意连续偶函数 $g(x)$,
\[
    \int_{-1}^{1}f(x)g(x)\,\mathrm{d}x=0.
\]
试证:$f(x)$ 是 $\lbrack -1,1\rbrack$ 上的奇函数.
\end{example}
\exampleblank

\begin{example}[特殊函数逼近]
设 $f\in R(a,b)$,证明:对任意 $\varepsilon>0$,存在函数 $g(x)$,使得
\[
    \int_a^b|f(x)-g(x)|\,\mathrm{d}x<\varepsilon,
\]
其中的 $g$ 是:
\begin{enumerate}[label=(\arabic*)]
    \item 阶梯函数;
    \item 折线函数;
    \item 连续函数;
    \item 连续可微函数.
\end{enumerate}
\end{example}
\exampleblank


\newpage

\section{积分估计与积分不等式}


\subsection{Newton-Leibniz 公式与上限变动积分}

\begin{example}
举例说明:可积不一定有原函数,原函数存在也不一定可积.
\end{example}
\exampleblank

\begin{theorem}[\markstar\ 原函数存在充分条件]
设 $f\in R(\lbrack a,b\rbrack)$,且在点 $x_0\in\lbrack a,b\rbrack$ 处连续,则其上限变动积分 $\displaystyle\int_a^x f(t)\,\mathrm{d}t$ 在点 $x=x_0$ 处可导,且其导数是 $f(x_0)$.
\end{theorem}

\begin{theorem}[\markstar]
若 $f\in R(\lbrack a,b\rbrack)$,则 $\displaystyle\int_a^x f(t)\,\mathrm{d}t$ 在 $\lbrack a,b\rbrack$ 上一致连续.
\end{theorem}

\begin{example}
设 $f\in C(\lbrack 0,1\rbrack)$ 且 $f(x)\geq0$,若有
\[
    f^2(x)\leq1+2\int_0^x f(t)\,\mathrm{d}t,
\]
则 $f(x)\leq1+x$.
\end{example}
\exampleblank

\begin{example}
\begin{enumerate}[label=(\arabic*)]
    \item 设 $f(x)$ 在 $\lbrack a,b\rbrack$ 上为正值可积函数,则存在 $\xi\in(a,b)$,使得
    \[
        \int_a^\xi f(x)\,\mathrm{d}x
        =\int_\xi^b f(x)\,\mathrm{d}x
        =\dfrac{1}{2}\int_a^b f(x)\,\mathrm{d}x;
    \]
    \item 设 $f\in C(\lbrack a,b\rbrack)$ 且 $\displaystyle\int_a^b f(x)\,\mathrm{d}x=0$,则存在 $\xi\in(a,b)$,使得
    \[
        \int_a^\xi f(x)\,\mathrm{d}x=f(\xi).
    \]
\end{enumerate}
\end{example}
\exampleblank

\begin{example}
\begin{enumerate}[label=(\arabic*)]
    \item 设 $f(x)$ 在 $\lbrack 0,1\rbrack$ 上可微,$0\leq f'(x)\leq1$ 且 $f(0)=0$,则
    \[
        \left(\int_0^1f(x)\,\mathrm{d}x\right)^2
        \geq\int_0^1f^3(x)\,\mathrm{d}x;
    \]
    \item 设 $f(x),g(x)$ 是 $\lbrack a,b\rbrack$ 上的递增连续函数,则
    \[
        \int_a^bf(x)\,\mathrm{d}x\int_a^bg(x)\,\mathrm{d}x
        \leq(b-a)\int_a^bf(x)g(x)\,\mathrm{d}x.
    \]
\end{enumerate}
\end{example}
\exampleblank

\begin{example}
设 $f\in C(\lbrack 0,1\rbrack)$,若有 $\displaystyle\int_0^x f(t)\,\mathrm{d}t\geq f(x)\geq0$,则 $f(x)\equiv0$.
\end{example}
\exampleblank

\begin{example}
设 $f\in C^{(1)}(\lbrack 0,a\rbrack)$,$f(0)=0$,则
\[
    \int_0^a|f(x)f'(x)|\,\mathrm{d}x
    \leq\dfrac{a}{2}\int_0^a[f'(x)]^2\,\mathrm{d}x.
\]
若条件改为 $f(a)=f(0)=0$ 或 $f\left(\dfrac{a}{2}\right)=0$,不等式右边系数该如何变化?
\end{example}
\exampleblank

\begin{example}[Gronwall 引理]
设 $f(x),g(x),\varphi(x)$ 均为 $\lbrack a,b\rbrack$ 上的连续函数,且 $g(x)$ 单调递增,$\varphi(x)\geq0$,且
\[
    f(x)\leq g(x)+\int_a^x\varphi(t)f(t)\,\mathrm{d}t.
\]
证明:对任意 $x\in\lbrack a,b\rbrack$,有
\[
    f(x)\leq g(x)e^{\int_a^x\varphi(s)\,\mathrm{d}s}.
\]
\end{example}
\exampleblank

\begin{example}
设 $f(x)\in C^1\lbrack 0,a\rbrack$,证明:
\[
    |f(0)|\leq\dfrac{1}{a}\int_0^a|f(x)|\,\mathrm{d}x
    +\int_0^a|f'(x)|\,\mathrm{d}x.
\]
\end{example}
\exampleblank

\begin{example}
设 $f(x)\in C^1\lbrack 0,a\rbrack$,证明:
\[
    \int_0^1|f(x)|\,\mathrm{d}x
    \leq\max\left\{\int_0^1|f'(x)|\,\mathrm{d}x+
    \left|\int_0^1f(x)\,\mathrm{d}x\right|\right\}.
\]
\end{example}
\exampleblank



\subsection{各种变形}

所谓各种变形即包括代数变形、换元(对称性)、分部积分、中值定理、Taylor 公式等.另外,被积函数放大、缩小,区间放大和缩小也是获得估计的重要方法.



\methodsection{(A) 放缩被积函数}

\begin{example}
证明:
\begin{enumerate}[label=(\arabic*)]
    \item
    \[
        \sqrt{2}<\int_{\sqrt{2}-1}^{\sqrt{2}}e^{-x^2}\,\mathrm{d}x<\sqrt{2};
    \]
    \item
    \[
        0<\dfrac{\pi}{2}-\int_0^{\frac{\pi}{2}}\dfrac{\sin x}{x}\,\mathrm{d}x
        <\dfrac{\pi^3}{144};
    \]
    \item
    \[
        \dfrac{2}{9}\pi^2
        \leq\int_{\frac{\pi}{6}}^{\frac{\pi}{2}}\dfrac{2x}{\sin x}\,\mathrm{d}x
        \leq\dfrac{4}{9}\pi^2.
    \]
\end{enumerate}
\end{example}
\exampleblank



\methodsection{(B) 换元法}

\begin{example}
证明:
\[
    \int_0^{\sqrt{2\pi}}\sin x^2\,\mathrm{d}x>0.
\]
\end{example}
\exampleblank

\begin{example}
$f\in C(\lbrack a,b\rbrack)$ 且严格递增,则
\[
    \int_a^bxf(x)\,\mathrm{d}x
    >\dfrac{a+b}{2}\int_a^bf(x)\,\mathrm{d}x.
\]
\end{example}
\exampleblank

\begin{example}
证明:
\[
    \int_0^{\frac{\pi}{2}}\dfrac{\cos x}{1+x^2}\,\mathrm{d}x
    >\int_0^{\frac{\pi}{2}}\dfrac{\sin x}{1+x^2}\,\mathrm{d}x.
\]
\end{example}
\exampleblank



\methodsection{(C) 分部积分}

\begin{proposition}[\markstar]
若 $f(x),g(x)$ 在 $\lbrack a,b\rbrack$ 上有 $2n$ 阶连续导数,且
\[
    f^{(i)}(a)=f^{(i)}(b)=g^{(i)}(a)=g^{(i)}(b)=0,
    \qquad i=0,1,2,\ldots,n-1,
\]
则
\[
    \int_a^bf^{(2n)}(x)g(x)\,\mathrm{d}x
    =\int_a^bf(x)g^{(2n)}(x)\,\mathrm{d}x.
\]
\end{proposition}

\begin{example}
设 $f$ 在 $\lbrack a,b\rbrack$ 上二阶可微,$f(a)=f(b)=0$,$|f''(x)|\leq M$,证明:
\[
    \left|\int_a^bf(x)\,\mathrm{d}x\right|
    \leq\dfrac{M}{2}(b-a)^3.
\]
\end{example}
\exampleblank

\begin{example}
设 $f(x)$ 在 $\lbrack a,b\rbrack$ 上有 $2n$ 阶连续导数,$|f^{(2n)}(x)|\leq M$,且
\[
    f^{(i)}(a)=f^{(i)}(b)=0,
    \qquad i=0,1,2,\ldots,n-1.
\]
试证:
\[
    \left|\int_a^bf(x)\,\mathrm{d}x\right|
    \leq\dfrac{(n!)^2M}{(2n)!(2n+1)!}(b-a)^{2n+1}.
\]
\end{example}
\exampleblank

\begin{example}
设 $f'(x)$ 在 $\lbrack 0,2\pi\rbrack$ 上连续,且 $f'(x)\geq0$,则对任意 $n\in\mathbb{Z}^+$,有
\[
    \left|\int_0^{2\pi}f(x)\sin nx\,\mathrm{d}x\right|
    \leq\dfrac{2[f(2\pi)-f(0)]}{n}.
\]
\end{example}
\exampleblank



\methodsection{(D) 中值定理}

\begin{example}
设函数 $f(x)$ 在 $\lbrack 2,4\rbrack$ 上有连续的一阶导数,且 $f(2)=f(4)=0$,试证:
\[
    \max\limits_{2\leq x\leq4}f'(x)
    \geq\int_2^4f(x)\,\mathrm{d}x.
\]
\end{example}
\exampleblank

\begin{example}
设 $f(x)$ 在 $\lbrack 0,1\rbrack$ 上有二阶连续导数,证明:
\[
    \max\limits_{x\in\lbrack 0,1\rbrack}|f'(x)|
    \leq|f(1)-f(0)|+\int_0^1|f''(x)|\,\mathrm{d}x.
\]
\end{example}
\exampleblank

\begin{example}
已知 $f(x)$ 在 $\lbrack 0,1\rbrack$ 上有二阶连续导数,且 $f(1)=f(0)=0$,证明:
\[
    \int_0^1|f'(x)|\,\mathrm{d}x
    \leq\int_0^1|f''(x)|\,\mathrm{d}x+2\int_0^1|f(x)|\,\mathrm{d}x.
\]
\end{example}
\exampleblank

\begin{example}
已知 $f(x)$ 在 $\lbrack 0,1\rbrack$ 上有直到二阶的连续导数,证明:
\[
    \int_0^1|f'(x)|\,\mathrm{d}x
    \leq\int_0^1|f''(x)|\,\mathrm{d}x+8\int_0^1|f(x)|\,\mathrm{d}x.
\]
\end{example}
\exampleblank

\begin{example}
$f(x)$ 在 $\lbrack 0,1\rbrack$ 上有连续二阶导数,$f(0)=f(1)=0$,$f(x)\neq0$(或 $f''(x)<0$),证明:
\[
    \int_0^1\left|\dfrac{f''(x)}{f(x)}\right|\,\mathrm{d}x\geq4.
\]
\end{example}
\exampleblank

\begin{example}
设 $f(x)\in C^2\lbrack a,b\rbrack$,且 $|f''(x)|\leq M$,证明:
\begin{enumerate}[label=(\arabic*)]
    \item 若 $f\left(\dfrac{a+b}{2}\right)=0$,则
    \[
        \left|\int_a^bf(x)\,\mathrm{d}x\right|
        \leq\dfrac{M}{24}(b-a)^3;
    \]
    \item 若 $f(a)=f(b)=0$,则
    \[
        \left|\int_a^bf(x)\,\mathrm{d}x\right|
        \leq\dfrac{M}{12}(b-a)^3.
    \]
\end{enumerate}
\end{example}
\exampleblank

\begin{example}
设 $f(x)\in C^2\lbrack a,b\rbrack$,满足
\[
    \left|\int_a^bf(x)\,\mathrm{d}x\right|
    <\int_a^b|f(x)|\,\mathrm{d}x,
\]
设
\[
    \max\limits_{x\in\lbrack a,b\rbrack}|f'(x)|=M_1,
    \qquad
    \max\limits_{x\in\lbrack a,b\rbrack}|f''(x)|=M_2.
\]
求证:
\[
    \left|\int_a^bf(x)\,\mathrm{d}x\right|
    \leq\dfrac{M_1}{2}(b-a)^2+\dfrac{M_2}{6}(b-a)^3.
\]
\end{example}
\exampleblank



\subsection{其它问题}

\begin{example}
设 $f(x)$ 在 $\lbrack 0,1\rbrack$ 上连续,
\[
    \int_0^1f(x)\,\mathrm{d}x=0,
    \quad
    \int_0^1xf(x)\,\mathrm{d}x=0,
    \quad\ldots,\quad
    \int_0^1x^{n-1}f(x)\,\mathrm{d}x=0,
\]
且
\[
    \int_0^1x^n f(x)\,\mathrm{d}x=1.
\]
求证:在 $\lbrack 0,1\rbrack$ 的某一部分上
\[
    |f(x)|\geq2^n(n+1).
\]
\end{example}
\exampleblank

\begin{example}
设 $f(x)$ 在 $\lbrack 0,\dfrac{\pi}{2}\rbrack$ 上连续,
\[
    \int_0^{\frac{\pi}{2}}f(x)\sin x\,\mathrm{d}x
    =\int_0^{\frac{\pi}{2}}f(x)\cos x\,\mathrm{d}x=0,
\]
证明:$f(x)$ 在 $\left(0,\dfrac{\pi}{2}\right)$ 内至少有两个零点.
\end{example}
\exampleblank

\begin{example}[Hadamard 定理]
设 $f(x)$ 是 $\lbrack a,b\rbrack$ 上连续凸函数,证明:对任意 $x_1,x_2\in\lbrack a,b\rbrack$ 且 $x_1<x_2$,有
\[
    f\left(\dfrac{x_1+x_2}{2}\right)
    \leq\dfrac{1}{x_2-x_1}\int_{x_1}^{x_2}f(t)\,\mathrm{d}t
    \leq\dfrac{f(x_1)+f(x_2)}{2}.
\]
\end{example}
\exampleblank

\begin{example}
设 $f(x)$ 是 $\lbrack 0,+\infty)$ 上的凸函数,求证:
\[
    F(x)=\dfrac{1}{x}\int_0^xf(t)\,\mathrm{d}t
\]
为 $(0,+\infty)$ 上的凸函数.
\end{example}
\exampleblank

\begin{example}[\markstar]
设 $f(x),g(x)$ 在区间 $\lbrack a,b\rbrack$ 上连续,$p(x)\geq0$,$\displaystyle\int_a^bp(x)\,\mathrm{d}x>0$,且
\[
    m\leq f(x)\leq M.
\]
$\varphi(x)$ 在 $\lbrack m,M\rbrack$ 上有定义,并有二阶导数,$\varphi''(x)>0$.求证:
\[
    \varphi\left(
    \dfrac{\displaystyle\int_a^bp(x)f(x)\,\mathrm{d}x}
    {\displaystyle\int_a^bp(x)\,\mathrm{d}x}
    \right)
    \leq
    \dfrac{\displaystyle\int_a^bp(x)\varphi(f(x))\,\mathrm{d}x}
    {\displaystyle\int_a^bp(x)\,\mathrm{d}x}.
\]
\end{example}
\exampleblank



\subsection{著名不等式}

\begin{theorem}[\marktwostar\ Schwarz 不等式]
设 $f,g\in R\lbrack a,b\rbrack$,则
\[
    \left(\int_a^bf(x)g(x)\,\mathrm{d}x\right)^2
    \leq\int_a^bf^2(x)\,\mathrm{d}x\int_a^bg^2(x)\,\mathrm{d}x.
\]
\end{theorem}

\begin{theorem}[\markstar\ Young 不等式]
设 $f$ 在 $\lbrack 0,+\infty)$ 上连续可微且严格单调递增,$f(0)=0$,$a,b>0$,则
\[
    ab\leq\int_0^af(x)\,\mathrm{d}x+\int_0^bg(y)\,\mathrm{d}y,
\]
其中 $g(y)$ 是 $f(x)$ 的反函数,而等号当且仅当 $b=f(a)$ 时成立.
\end{theorem}

\begin{theorem}[\markstar\ Holder 不等式]
设 $f,g\in R\lbrack a,b\rbrack$,$p,q$ 为满足 $\dfrac{1}{p}+\dfrac{1}{q}=1$ 的一对正实数,则
\[
    \int_a^b|f(x)g(x)|\,\mathrm{d}x
    \leq
    \left(\int_a^b|f(x)|^p\,\mathrm{d}x\right)^{\frac{1}{p}}
    \left(\int_a^b|g(x)|^q\,\mathrm{d}x\right)^{\frac{1}{q}}.
\]
\end{theorem}

\begin{theorem}[\markstar\ Minkowski 不等式]
设 $f,g\in R\lbrack a,b\rbrack$,$1\leq p\leq+\infty$,则
\[
    \left(\int_a^b|f(x)+g(x)|^p\,\mathrm{d}x\right)^{\frac{1}{p}}
    \leq
    \left(\int_a^b|f(x)|^p\,\mathrm{d}x\right)^{\frac{1}{p}}
    +\left(\int_a^b|g(x)|^p\,\mathrm{d}x\right)^{\frac{1}{p}}.
\]
\end{theorem}

\begin{example}
设 $f\in C\lbrack 0,1\rbrack$,$0\leq f(x)<1$,则
\[
    \int_0^1\dfrac{f(x)}{1-f(x)}\,\mathrm{d}x
    \geq
    \dfrac{\displaystyle\int_0^1\sqrt{f(x)}\,\mathrm{d}x}
    {1-\displaystyle\int_0^1f(x)\,\mathrm{d}x}.
\]
\end{example}
\exampleblank

\begin{example}
\begin{enumerate}[label=(\arabic*)]
    \item 设 $f\in R\lbrack a,b\rbrack$,则
    \[
        \left(\int_a^bf(x)\sin x\,\mathrm{d}x\right)^2
        +\left(\int_a^bf(x)\cos x\,\mathrm{d}x\right)^2
        \leq(b-a)\int_a^bf^2(x)\,\mathrm{d}x.
    \]
    \item 设 $f\in R\lbrack a,b\rbrack$ 且 $f(x)>0$,则
    \[
        \left(\int_a^bf(x)\sin x\,\mathrm{d}x\right)^2
        +\left(\int_a^bf(x)\cos x\,\mathrm{d}x\right)^2
        \leq\left(\int_a^bf(x)\,\mathrm{d}x\right)^2.
    \]
\end{enumerate}
\end{example}
\exampleblank

\begin{example}
设 $f\in R\lbrack a,b\rbrack$,若 $0<m\leq f(x)\leq M$,则:
\begin{enumerate}[label=(\arabic*)]
    \item
    \[
        \left(\int_a^bf(x)\,\mathrm{d}x\right)
        \left(\int_a^b\dfrac{\mathrm{d}x}{f(x)}\right)
        \leq\dfrac{(m+M)^2}{4mM}(b-a)^2;
    \]
    \item
    \[
        2\sqrt{\dfrac{m}{M}}(b-a)
        \leq\dfrac{1}{M}\int_a^bf(x)\,\mathrm{d}x
        +m\int_a^b\dfrac{\mathrm{d}x}{f(x)}.
    \]
\end{enumerate}
\end{example}
\exampleblank

\begin{example}
设 $a>0$,证明:
\[
    \int_0^\pi x^a\sin x\,\mathrm{d}x
    \int_0^{\frac{\pi}{2}}a^{-\cos x}\,\mathrm{d}x
    \geq\dfrac{\pi^3}{4}.
\]
\end{example}
\exampleblank

\begin{example}
设 $f(x)\in C^1\lbrack a,b\rbrack$,$f(a)=f(b)=0$,$\displaystyle\int_a^bf^2(x)\,\mathrm{d}x=1$,证明:
\[
    \int_a^b[f'(x)]^2\,\mathrm{d}x
    \int_a^b[xf(x)]^2\,\mathrm{d}x
    \geq\dfrac{1}{4}.
\]
\end{example}
\exampleblank

\begin{example}
设 $f(x)\in C^1\lbrack a,b\rbrack$,存在 $c\in\lbrack a,b\rbrack$,使得 $f(c)=0$.记
\[
    M=\max\limits_{x\in\lbrack a,b\rbrack}|f(x)|,
\]
证明:
\[
    M^2\leq(b-a)\int_a^b[f'(x)]^2\,\mathrm{d}x.
\]
\end{example}
\exampleblank

\begin{example}
设 $f(x)$ 在 $(0,1)$ 上可积,且 $\displaystyle\int_0^1xf(x)\,\mathrm{d}x=0$,证明:
\[
    \int_0^1f^2(x)\,\mathrm{d}x
    \geq4\left(\int_0^1f(x)\,\mathrm{d}x\right)^2.
\]
\end{example}
\exampleblank


\newpage

\section{积分中值定理}

\begin{theorem}[\markstar\ 第一积分中值定理]
设 $f(x)$ 是 $\lbrack a,b\rbrack$ 上的连续函数,则:
\begin{enumerate}[label=(\arabic*)]
    \item (强化形式)存在 $\xi\in(a,b)$,使得
    \[
        \int_a^bf(x)\,\mathrm{d}x=f(\xi)(b-a).
    \]
    \item (推广形式)设 $g(x)$ 在 $\lbrack a,b\rbrack$ 上可积且不变号,则存在 $\xi\in(a,b)$,使得
    \[
        \int_a^bf(x)g(x)\,\mathrm{d}x
        =f(\xi)\int_a^bg(x)\,\mathrm{d}x.
    \]
\end{enumerate}
\end{theorem}

\begin{remark}
由证明过程可得,替代连续性,第一积分中值定理对具有介值性的可积函数仍然成立.
\end{remark}

\begin{example}
设 $f(x)$ 在 $(0,1)$ 上可微,且满足
\[
    f(1)-2\int_0^1xf(x)\,\mathrm{d}x=0.
\]
求证:在 $(0,1)$ 内至少有一点 $\xi$,使得
\[
    f'(\xi)=-\dfrac{f(\xi)}{3}.
\]
\end{example}
\exampleblank

\begin{example}
设 $f(x)$ 在 $\lbrack 0,1\rbrack$ 上连续,在 $(0,1)$ 内三次可微,并且满足条件:
\begin{enumerate}[label=(\arabic*)]
    \item $\displaystyle\int_0^{\frac{1}{4}}f(t)\,\mathrm{d}t=\dfrac{1}{4}f(0)$;
    \item $\displaystyle\int_{\frac{1}{4}}^1f(t)\,\mathrm{d}t=\dfrac{3}{4}f(1)$.
\end{enumerate}
试证:若 $f(0)=f(1)$,则存在 $\xi\in(0,1)$,使得 $f'''(\xi)=0$.
\end{example}
\exampleblank

\begin{example}
$\lbrack a,b\rbrack$ 上的连续函数序列 $\varphi_1,\varphi_2,\ldots,\varphi_n,\ldots$ 满足
\[
    \int_a^b\varphi_n(x)\,\mathrm{d}x=1.
\]
证明:存在自然数 $N$ 及定数 $c_1,c_2,\ldots,c_N$,使得
\[
    \sum\limits_{k=1}^N c_k^2=1,
    \qquad
    \max\limits_{x\in\lbrack a,b\rbrack}\left|\sum\limits_{k=1}^Nc_k\varphi_k(x)\right|>100.
\]
\end{example}
\exampleblank

\begin{example}
求
\[
    \lim\limits_{x\to+\infty}
    \int_x^{x+2}t\left(\sin\dfrac{3}{t}\right)f(t)\,\mathrm{d}t,
\]
其中 $f(x)$ 可微,且 $\displaystyle\lim\limits_{t\to+\infty}f(t)=1$.
\end{example}
\exampleblank

\begin{theorem}[\markstar\ 第二积分中值定理]
\begin{enumerate}[label=(\arabic*)]
    \item 在 $\lbrack a,b\rbrack$ 上,若 $f(x)$ 非负单调递减,$g(x)$ 可积,则存在 $\xi\in\lbrack a,b\rbrack$,使得
    \[
        I=\int_a^bf(x)g(x)\,\mathrm{d}x
        =f(a)\int_a^\xi g(x)\,\mathrm{d}x.
    \]
    \item 在 $\lbrack a,b\rbrack$ 上,若 $f(x)$ 非负单调递增,$g(x)$ 可积,则存在 $\xi\in\lbrack a,b\rbrack$,使得
    \[
        I=\int_a^bf(x)g(x)\,\mathrm{d}x
        =f(b)\int_\xi^bg(x)\,\mathrm{d}x.
    \]
    \item 在 $\lbrack a,b\rbrack$ 上,若 $f(x)$ 单调,$g(x)$ 可积,则存在 $\xi\in\lbrack a,b\rbrack$,使得
    \[
        I=\int_a^bf(x)g(x)\,\mathrm{d}x
        =f(a)\int_a^\xi g(x)\,\mathrm{d}x
        +f(b)\int_\xi^bg(x)\,\mathrm{d}x.
    \]
\end{enumerate}
\end{theorem}

\begin{remark}
因此,若见积分问题条件中有单调性,可优先考虑第二积分中值定理.且由证明过程可知,第一式中的 $f(a)$ 可改写为 $f(a+0)$ 或比 $f(a+0)$ 大的任意常数 $A$,即存在 $\xi\in\lbrack a,b\rbrack$,使得
\[
    I=\int_a^bf(x)g(x)\,\mathrm{d}x
    =A\int_a^\xi g(x)\,\mathrm{d}x.
\]
其中,$\xi$ 依赖于 $A$ 的选取.同理,第二式中的 $f(b)$ 可改写为 $f(b-0)$ 或比 $f(b-0)$ 大的任一常数 $B$.
\end{remark}

\begin{example}
证明:若 $f(x)$ 在 $\lbrack 0,1\rbrack$ 上严格单调递减,则:
\begin{enumerate}[label=(\arabic*)]
    \item 存在 $\theta\in(0,1)$,使得
    \[
        \int_0^1f(x)\,\mathrm{d}x=\theta f(0)+(1-\theta)f(1);
    \]
    \item 对任意 $c>f(0)$,存在 $\theta\in(0,1)$,使得
    \[
        \int_0^1f(x)\,\mathrm{d}x=\theta c+(1-\theta)f(1).
    \]
\end{enumerate}
\end{example}
\exampleblank

\begin{example}
设 $f:\lbrack a,b\rbrack\to\lbrack 0,1\rbrack$ 可导,$f'(x)$ 在 $\lbrack a,b\rbrack$ 上单调递减,$f'(b)>0$,求证:
\[
    \int_a^b\cos f(x)\,\mathrm{d}x\leq\dfrac{2}{f'(b)}.
\]
\end{example}
\exampleblank

\begin{example}
\begin{enumerate}[label=(\arabic*)]
    \item 设 $x>0,c>0$,证明:
    \[
        \left|\int_x^{x-c}\sin t^2\,\mathrm{d}t\right|<\dfrac{1}{x};
    \]
    \item 设 $e^2<a<b$,证明:
    \[
        \int_a^b\dfrac{\mathrm{d}x}{\ln x}<\dfrac{2b}{\ln b};
    \]
    \item 设 $0<a<b$,证明:
    \[
        \left|\int_a^b\dfrac{\sin x}{x}\,\mathrm{d}x\right|\leq\dfrac{2}{a}.
    \]
\end{enumerate}
\end{example}
\exampleblank

\begin{example}
设 $f(x)$ 是 $\lbrack -\pi,\pi\rbrack$ 上的凸函数,$f'(x)$ 有界,求证:
\[
    a_{2n}=\dfrac{1}{\pi}\int_{-\pi}^\pi f(x)\cos2nx\,\mathrm{d}x\geq0,
\]
\[
    a_{2n+1}=\dfrac{1}{\pi}\int_{-\pi}^\pi f(x)\cos(2n+1)x\,\mathrm{d}x\leq0.
\]
\end{example}
\exampleblank

\begin{example}
设 $f(x)\geq0$ 在 $\lbrack -\pi,\pi\rbrack$ 上单调递减,证明:
\[
    a_{2n}=\dfrac{1}{\pi}\int_{-\pi}^\pi f(x)\sin2nx\,\mathrm{d}x\geq0,
\]
\[
    a_{2n+1}=\dfrac{1}{\pi}\int_{-\pi}^\pi f(x)\sin(2n+1)x\,\mathrm{d}x\leq0.
\]
\end{example}
\exampleblank

\begin{example}
证明下列命题:
\begin{enumerate}[label=(\arabic*)]
    \item
    \[
        \int_0^1\dfrac{\mathrm{d}x}{1+x^n}
        =1-\dfrac{k^2}{n}+o(1),
        \qquad n\to\infty;
    \]
    \item
    \[
        \lim\limits_{x\to0}\dfrac{1}{x}\int_x^{2x}\sin\dfrac{1}{t}\,\mathrm{d}t=0;
    \]
    \item 当 $\alpha<2$ 时,
    \[
        \lim\limits_{x\to0^+}\dfrac{1}{x^\alpha}\int_0^x\sin\dfrac{1}{t}\,\mathrm{d}t=0.
    \]
\end{enumerate}
\end{example}
\exampleblank

\end{document}
